\subsection{Attribute Program Synthesis}
\label{sec:synthesis}
TEPL uses symbolic axis groups and mappings to turn the construction in
Section~\ref{sec:mathematical-synthesis} into reusable attribute programs.


\subsubsection{Match Source and Replacement Computations}
\label{sec:normalization}
The compiler first identifies which source coordinates correspond to each
RHS operand read. It normalizes the source computation and the RHS
\emph{arithmetic skeleton} into sums of products, similar to SPORES's
R-polynomial representation~\cite[Section~2.3, Definition~2.4]{spores},
then matches reads of the same input tensors. These matches associate
RHS reads with source coordinates before the RHS axis fields are known.
For reassociation, both computations contain reads of $X,Y,Z$; their
completed reduction structure and equivalence are checked after construction.

The replacement intermediate $U=\operatorname{dot}(Y,Z)$ uses the matched
reads of $Y,Z$, while the final dot combines $X$ and $U$.
Appendix~\ref{extended:sec:normalization} details normalization and read
matching.


\subsubsection{Construct Attribute Programs}
\label{sec:keep-axes}
\label{sec:attribute-construction}
\label{sec:emit-programs}
Using Section~\ref{sec:mathematical-synthesis}'s rules, the compiler selects
axis groups for each replacement operation. Symbolic mappings convert
these groups into operand positions: $\pos(\mathsf I,\vec G)$ locates
the ordered group $\vec G$ in interface $\mathsf I$.
For $U$, both contraction lists select $\vec J_{\mathrm{sel}}$:
\begin{equation}
 k_\ell^U=\pos(\mathsf{I}_Y,\vec J_{\mathrm{sel}}),\qquad
 k_r^U=\pos(\mathsf{I}_Z,\vec J_{\mathrm{sel}}).
 \label{eq:symbolic-contraction-queries}
\end{equation}
Both queries select the same group order, preserving operand pairings.

\sidefigure{figures/Constructing Dot Attributes from Axis Mappings.pdf}{Paired position queries construct dot attributes; the completed intermediate supplies mappings for the final dot.}{fig:attribute-construction}{Paired queries on Y and Z construct the inner dot, an interface bridge
establishes U's layout, and paired queries on X and U construct the final dot.}


Figure~\ref{fig:attribute-construction} shows construction of $U$ followed
by the final dot. After constructing $U$, the compiler records its ordered
interface $\mathsf I_U$. In the figure's unbatched example, the final dot
queries this interface for its contraction group
$\vec G=\vec I_{\mathrm{sel}}\catseq\vec K$.
The interface preserves surviving operand order, including interleaved groups.
Appendix~\ref{extended:sec:attribute-construction} gives the full derivation.

The compiler resolves position queries into expressions over source
attributes and shapes. Queries involving a replacement intermediate
substitute its constructed interface until only source metadata remains.
The resulting programs contain no unresolved axis fields or logical-group
queries. They are generated once per rule and evaluated on each match,
allowing attribute lists to vary with tensor rank.
Appendix~\ref{extended:sec:emit-programs} details program generation
and verification for the reassociation example.

TEPL verifies the completed rewrite under source legality and the
declared \code{where} condition, as described in
Section~\ref{sec:condition-suggestions}.
